\documentclass[11pt,a4paper]{article}
\usepackage[T1]{fontenc}
\usepackage[utf8]{inputenc}
\usepackage{lmodern,microtype}
\usepackage[margin=.85in]{geometry}
\usepackage{amsmath,amssymb,amsthm,mathtools,aliascnt}
\usepackage{booktabs,tabularx,array,enumitem,needspace,float}
\usepackage[section]{placeins}
\usepackage{xcolor,tikz,mdframed}
\usetikzlibrary{arrows.meta,positioning,calc}
\usepackage[hidelinks]{hyperref}
\usepackage[nameinlink,noabbrev]{cleveref}
\usepackage{fancyhdr}
\hypersetup{pdftitle={Sync HotStuff RR v4},pdfsubject={Stake-weighted regional sessions and cross-region bootstrap}}
\setlength{\headheight}{14pt}
\pagestyle{fancy}\fancyhf{}
\lhead{\small Sync HotStuff RR v4}\rhead{\small L1-authorized regional sessions}
\cfoot{\thepage}
\setlength{\parindent}{0pt}\setlength{\parskip}{4pt}
\setlength{\emergencystretch}{2em}
\setlist{nosep,leftmargin=*}
\newcolumntype{Y}{>{\raggedright\arraybackslash}X}
\definecolor{rrRed}{HTML}{A51C30}
\definecolor{rrBlue}{HTML}{2563EB}
\definecolor{rrTeal}{HTML}{0F766E}
\newcommand{\RR}[1]{{\color{rrRed}#1}}
\newcommand{\QC}{\mathcal C}
\newcommand{\FC}{\mathcal F}
\newcommand{\Hash}{\mathsf H}
\newcommand{\rank}{\operatorname{rank}}
\newcommand{\wt}{\operatorname{wt}}
\newcommand{\tagg}[1]{\ensuremath{\mathsf{#1}}}
\newtheorem{definition}{Definition}[section]
\newaliascnt{assumption}{definition}
\newtheorem{assumption}[assumption]{Assumption}
\aliascntresetthe{assumption}
\crefname{assumption}{assumption}{assumptions}
\Crefname{assumption}{Assumption}{Assumptions}
\newcounter{protocol}
\newenvironment{protocol}[2]{\refstepcounter{protocol}\label{#1}
\begin{mdframed}[nobreak=true,linewidth=.6pt,skipabove=8pt,skipbelow=8pt,
innertopmargin=7pt,innerbottommargin=7pt]\small
\textbf{Protocol \theprotocol. #2}\par\smallskip}{\end{mdframed}}
\crefname{protocol}{Protocol}{Protocols}
\Crefname{protocol}{Protocol}{Protocols}
\title{Sync HotStuff RR v4\\\large Weighted Voting, Restricted Drain, and Regional Bootstrap}
\author{}\date{12 October 2026}
\begin{document}
\maketitle
\begin{abstract}
L1 fixes regional stake snapshots, leader assignments, application-block
quotas, and scheduled status updates. A global view $(r,g,t)$ identifies
a regional session, recovery generation, and leader term. Within a region,
terminal certificates rotate leaders while proposal timers continue.
At the regional quota, the terminal DATA block completes on its own timer;
one DRAIN block is permitted only in a recovery generation whose opening
already contains the full quota. Prospective successors continuously fetch
the current stream and finalized L1 control state. A regular or activated
forced-handoff event invokes one bootstrap handler with a finalized DATA
anchor. This draft specifies the model and protocols; security and
complexity analyses are deferred.
\end{abstract}

\section{Model and L1 control state}
\label{sec:model}
L1 is the public source of sequencing authority. It determines who may
vote and lead, how many DATA blocks a region may add, which regions are
eligible, and when an explicit forced rotation becomes effective.
Regional Sync HotStuff~\cite{sync} orders and commits the blocks.
Cross-region infrastructure distributes their contents, certificates, and
L1 proofs; it does not elect a leader by reporting a tip.

\subsection{Objects read from finalized L1}
\label{sec:l1}
A regional session $r\ge0$ is one authorized production instance of a
physical region $i_r$. Returning to the same physical region creates a
new session. Nodes read the following objects, with L1 inclusion and
finality proofs. A schedule may be represented by a deterministic rule
rather than an infinite list.
An authority may likewise be the instantiation of finalized L1 rules in
an authenticated session context: the boundary determines $i_r$, and the
configured rule supplies its leaders and quota. This does not require
posting a new leader-assignment transaction at every regular boundary.

\begin{table}[htbp]\centering\small
\caption{L1 control objects and their use by replicas.}
\label{tab:l1}
\begin{tabularx}{\linewidth}{@{}p{.25\linewidth}Y@{}}
\toprule
Object & Authenticated contents and meaning\\
\midrule
$\tagg{CFG}$ & Configuration identifier; regional member sets $V_i$, public
keys, stake weights $\omega_{i,v}$, fault bounds $F_i$, delay bounds
$\Delta_i$, local quota $K$, and cyclic region order $\mathcal L_R$.\\
$\tagg{AUTH}_{r,g}$ & The session's configuration, region $i_r$, fixed
DATA quota $TB_r\ge1$, and leader rule
$L(r,g,t)=(i_r,\ell_{r,g,t})$. A recovery authority preserves
$i_r$, the voting snapshot, and $TB_r$.\\
$\widehat E_j$ & An ordered multisignature entry containing a regional
status map $\chi_j$, future L1 activation height $a_j$, and either no
force tag or $\tagg{FORCE}(r_j,h_j^*)$. The tag names the retiring
session and its last available DATA height, not a block hash.\\
L1 proof at $\lambda$ & Canonical L1 block hash and authenticated registry
state, including the complete installed entry prefix. This distinguishes
an announced entry from one effective at $\lambda$.\\
\bottomrule
\end{tabularx}
\end{table}

Write the status entry and its signed envelope as
\[
 E_j=(a_j,\chi_j,\phi_j),\qquad
 \phi_j\in\{\bot,\tagg{FORCE}(r_j,h_j^*)\},\qquad
 \widehat E_j=(D_C,j,p_j,E_j,\Sigma_j^C),
\]
where $p_j$ commits to the preceding entry and $D_C$ binds the chain,
registry, and multisignature configuration. The map
$\chi_j:\mathcal R\to\{\tagg{live},\tagg{off}\}$ determines eligibility.
A force tag also marks its target region off. A passive entry changes
eligibility without itself replacing the active session. The multisignature
authorization policy is separate from regional vote thresholds.
The initial configuration installs $E_0=(a_0,\chi_0,\bot)$ and genesis
with a synthetic finality receipt. At the initial control height, the
first live region in $\mathcal L_R$ receives the synthetic bootstrap
boundary for session $0$; if none is live, initialization waits for an
eligible effective entry. All control heights below are at least $a_0$.

The registry enforces consecutive versions, the predecessor link, and
strictly increasing activation heights. An entry installed at L1 height
$u_j$ must satisfy $a_j\ge u_j+w_{\min}$ for a configured notice distance
$w_{\min}>0$. Finalized entries and authorities are immutable. At a
finalized L1 control height $\lambda$, the effective map is
\[
 j(\lambda)=\max\{j:a_j\le\lambda\},\qquad
 \chi_{\rm eff}(\lambda)=\chi_{j(\lambda)}.
\]
Installation alone does not activate a force tag. Increasing an L1 height
by several blocks is not by itself a wall-clock preparation bound.

\begin{assumption}[Finalized authority and ordered control]
\label{ass:authority}
L1 supplies one finalized control history and progresses independently
of regional commitment. For every admitted $(r,g)$, the matching authority
is available at a finite common time $d(r,g)$ to all correct members of
$V_{i_r}$, including frozen members. Control advances and sealed regional
handoffs have one authenticated order. Their witnesses identify the active
session, effective status prefix, and unique successor of each closed
boundary. A sealed handoff cannot be replaced by a later control update.
\end{assumption}
The control-order interface is inherited from the cross-region service.
It includes advances while block production is silent. No additional L1
transaction is required by the regional steady-state protocol or by each
regular handoff. Authenticating this common order is a service obligation;
local receipt order is not a substitute.

\subsection{Stake, communication, and recovery services}
For region $i$, let $W_i=\sum_{v\in V_i}\omega_{i,v}$ with
$\omega_{i,v}>0$. Byzantine stake is at most $F_i<W_i/2$.
Every regional certificate in this draft requires distinct signers of
total weight strictly greater than $F_i$. Equivalently, for a fault-share
bound $f_i=F_i/W_i<1/2$, its weight exceeds $f_iW_i$.
Membership and weights are fixed throughout a session, including local
recoveries. Slashing or a passive status entry does not change its
denominator. A later session verifies old evidence against the old snapshot.

\begin{assumption}[Authenticated regional synchrony and global clock]
\label{ass:network}
While a regional session operates under the normal model, every
correct-to-correct message, including its validation data and processing,
arrives within $\Delta_i$; actual delay may be smaller. Replicas share a
global clock, retain their protocol state, and have unforgeable signatures
and binding hashes. Correct members of a prospective successor also
receive and validate its bootstrap bundles while awaiting activation.
\end{assumption}
The fault bound and synchrony condition are per region, not merely bounds
on the aggregate network. Cross-region catch-up and L1 dissemination have
their own service delays. They are not included silently in $\Delta_i$.

\begin{assumption}[Cross-region retrieval and bootstrap readiness]
\label{ass:infra}
Prospective successors continuously receive or fetch leading proposals,
QCs, finality receipts, ancestry, and finalized L1 control proofs through
the cross-region infrastructure. It provides leader push and bounded-range
poll--fetch repair from retained holders. For an admitted successor, the
complete bootstrap bundle and required validation state reach its correct
members with entry skew at most $\Delta_i$. Each first recipient relays
the bundle within that region. Missing cross-region data are recovered
before this admission condition is met.
\end{assumption}
An unconfirmed tip is a retrieval hint. Neither a fetch timeout nor an
unsigned frontier grants production authority.

\begin{assumption}[Multisignature last-available height and forced closure]
\label{ass:cut}
A force tag for session $r$ identifies a last available DATA height $h_*$.
Through that height, regional BFT operated correctly and its DATA prefix
was broadcast and committed. The corresponding block, finality receipt,
and missing ancestry are recoverable outside the retiring region.
Its accepted closure witness additionally guarantees that this prefix
contains every completed honest DATA commitment, delivered final prefix,
and prefix for which a final receipt can already be assembled under the
affected authority. Old authority is closed before successor admission:
it cannot create a later accepted continuation outside the retained prefix
and new authorized path. Closure reconciles pending work and concurrent
handoffs without revoking a completed DATA commitment.
\end{assumption}
These are obligations of the multisignature/control service, rather than
consequences of elapsed L1 waiting. They permit recovery without a fresh
response from the failed region. The regional engine alone does not
establish them after an entire-region outage or partition. A higher valid
final prefix contradicts an undersized cut; it is not discarded because
its receipt arrived late. A status signature specifies no anchor hash;
regional finality evidence determines the block at the selected height.

\paragraph{Participation and retained work.}
Correct leaders may propose empty DATA blocks. Persistently admissible
client requests and their ordering state remain pending until committed,
including across sessions. Correct members retain and serve validation
data. Leader fairness and eventual exclusion of equivocating leaders are
the additional progress premises used in v3.1; changing stake identities
over time requires a corresponding scheduling and accountability policy.
This draft does not claim unconditional progress from economic stake alone.

\section{Views, blocks, and admission}
\label{sec:definitions}
A global view is $(r,g,t)$, ordered lexicographically. Here $r$ counts
regional sessions, $g$ counts recoveries within one session, and $t$ counts
ordinary leader terms within one generation. Each newly opened region
starts at $(r,0,0)$.

\begin{table}[htbp]\centering\small
\caption{View changes and their admission events.}
\label{tab:views}
\begin{tabularx}{\linewidth}{@{}p{.24\linewidth}p{.27\linewidth}Y@{}}
\toprule
Event & Target & Admission and timer treatment\\
\midrule
In-region normal rotation & $(r,g,t+1)$ & A term-terminal DATA QC below
the regional quota. Existing proposal timers continue.\\
In-region Quit & $(r,g+1,0)$ & A complete quit proof, the recovery wait,
and the next authority. Unfinished old-generation timers are canceled.\\
Regional quota completion & $(r+1,0,0)$ & A finalized terminal DATA anchor
and a sealed regular-handoff witness naming the successor.\\
Activated force & $(r+1,0,0)$ & Finalized L1 activation, the targeted
closure witness, and the finalized DATA anchor at the specified height.\\
\bottomrule
\end{tabularx}
\end{table}
A DRAIN QC causes no view change. A status update without an applicable
force tag does not increment $r$, $g$, or $t$. Local recovery never resets
the regional quota.

\subsection{Two positions and the regional quota}
\label{sec:positions}
The unchanged DATA block $A_r$ is the session's cross-region anchor;
$A_0$ is genesis. Let $h(B)$ be global DATA height and $k_r(B)$ the
physical position in session $r$. Assign $k_r(A_r)=0$ for this session's
arithmetic, without rewriting the anchor's original signed fields.
A fresh physical block has the form
\[
 B=(r,k,h,\mathit{type},x,\Hash(B_{\rm parent}),\rho),
 \qquad \mathit{type}\in\{\tagg{DATA},\tagg{DRAIN}\},
\]
where $k=k_r(B_{\rm parent})+1$. DATA contains the next application batch
and has $h=h(B_{\rm parent})+1$. DRAIN has empty application contents
and $h=h(B_{\rm parent})$. Define
\[
 H_r=h(A_r)+TB_r.
\]
DATA ancestry may contain at most $TB_r$ fresh DATA blocks, including
uncommitted ones. An ANCHOR re-certification, a retransmission, or a DRAIN
does not consume this quota. The terminal DATA block $T_r$ has height $H_r$.

Each DATA block binds a finalized L1 reference $\rho$ with height
$\lambda(B)$, canonical hash, and registry proof. These heights do not
decrease along DATA ancestry; the first fresh DATA reference is at least
the session's opening control height. Fresh admission verifies authority
against the complete control prefix effective at that reference.
DRAIN preserves its terminal DATA ancestor's reference and cannot delay
an effective forced closure. Historical block validity is not rewritten
using a receiver's newer L1 head. Closure separately prevents new accepted
work under retired authority, including work carrying an old reference.

\paragraph{Ordinary and shortened terms.}
Generation $(r,g)$ has an immutable opening $A_{r,g}$; $A_{r,0}=A_r$.
If $h(A_{r,g})<H_r$, term $t$ has
\[
 m_{r,g,t}=\min\{K,H_r-h(A_{r,g})-tK\}>0
\]
fresh DATA positions $j=0,\ldots,m_{r,g,t}-1$. Their positions are
\[
 k=k_r(A_{r,g})+tK+j+1,\qquad
 h=h(A_{r,g})+tK+j+1.
\]
There is no next term once its terminal DATA reaches $H_r$.
For example, $K=4$ and $TB_r=10$ give term lengths $4,4,2$ in an
uninterrupted generation starting at $A_r$.

\paragraph{The one permitted DRAIN position.}
\label{sec:drain}
If $g>0$ and $h(A_{r,g})=H_r$, this is a \emph{drain generation}.
Its only fresh position is $(t,j)=(0,0)$: one DRAIN extending $A_{r,g}$
with its current ANCHOR QC. The opening may itself be an older DRAIN.
There is no second fresh position or ordinary next term in this generation.
A generation that reaches $H_r$ by proposing DATA cannot append DRAIN;
it waits for commitment or recovers after Quit. Thus DRAIN is confined
to recovery from a full-quota opening. Repeated interrupted recoveries may
each add one DRAIN; there is no fixed bound across arbitrarily many such
generations. Whether some other replica has already committed the terminal
block is not an admission test.

\paragraph{Application finality and control history.}
Committing a physical block commits its ancestors. Across regions, the
replicated application log is the DATA projection: its committed prefixes
must be comparable and never revoked. The next session's first DATA block
names $\Hash(T_r)$, or the forced DATA anchor, as its parent. Old DRAIN
descendants remain regional control history and validation evidence; they
are not additional application positions or alternative cross-region
anchors. Within a session, physical locks and certificates retain the full
DATA/DRAIN ancestry.

\subsection{Certificates, receipts, and local processing}
Regional protocol signatures other than FINAL bind the chain, voting configuration, L1 authority,
session and region, generation opening, message type, view, physical
position, DATA height, and block hash. PROPOSE also binds $j$ and block
type. A proposal QC $\QC_{r,g,t}(B)$ contains matching votes of weight
$>F_{i_r}$. ANCHOR votes are separately tagged and re-certify an unchanged
opening. Correct replicas vote at most once per fresh $(r,g,t,j)$ and
once for an ANCHOR per $(r,g)$. DATA and DRAIN share the fresh-proposal
vote rule and differ in their signed block type.

Within a session, certificate rank is $(g,t,k_r(B))$. Equal-rank choices
use a fixed total tie-break on certified content, independently of signer
subsets or encodings. An opening's identity includes its hash and positions.
The verified cross-region anchor has a distinguished base justification
$Q_r^0$ below this session's ordinary QCs. It remains an admissible
recovery fallback if no ordinary QC has formed. An unrelated region's high-ranked QC
cannot replace that anchor.

\paragraph{Finality receipts.}
After committing a DATA block $B$, directly or through a descendant,
a correct replica sends
\[
 \langle\tagg{FINAL},\mathit{cfg}_r,r,h(B),\Hash(B)\rangle_v
\]
to its peers and the cross-region infrastructure. Matching signatures of
weight $>F_{i_r}$ form $\FC_r(B)$. A receipt includes the block and required
ancestry, or resolves them from a validated cache. It certifies commitment,
whereas a proposal QC alone does not. If a DRAIN timer commits $T_r$,
the exported terminal receipt names $T_r$, not that DRAIN.
A valid finality receipt may be adopted to commit its DATA prefix locally.
The FINAL signing domain binds the chain and the displayed session-stable
fields, independently of the generation or term in which commitment
occurs. Endorsements of the same DATA block therefore combine across
recoveries; the block hash preserves its original signed context.

\paragraph{Complete validation and evidence priority.}
For $g>0$, NEWVIEW is justified by a valid QC from an older generation
of the same session, or its base justification $Q_r^0$, and its quit-entry
proof. For $g=0$, it must use the
cross-region handler's exact anchor and bootstrap proof. Its ANCHOR QC
admits the first fresh proposal. Term $t>0$ additionally requires all
preceding term-terminal QCs in that generation, each below $H_r$.
Genesis is the sole synthetic initial anchor.

Validation checks signatures, authority, opening and term admission,
type and quota, both positions, parent QCs, hash links, and application
validity throughout the required ancestry. Every relevant ancestor includes
its signed proposal and block data. Correct transmissions provide complete
bundles; cached references must resolve within the complete-message bound.
Inspect signed material for independently verifiable quit evidence before
rejecting an enclosing message or discarding a duplicate. Compare it with
retained histories. At an equal timestamp, process received messages,
including evidence and complete proposals, before timer callbacks.

\paragraph{One proposal-processing rule.}
Every previously unvoted, admitted current-generation DATA or DRAIN
proposal invokes Vote, including attached history, earlier-term DATA,
and blocks already committed as ancestors. Process ancestors before
children and before terminal-QC entry. Deduplicate by
$(r,g,t,j,\Hash(B))$ after evidence inspection. Old-generation proposals
are retained but receive no new votes or timers; future-generation
proposals wait for entry. A certificate or an ANCHOR vote starts no timer.

\paragraph{Active, frozen, and complete sessions.}
New proposals, votes, and normal term transitions require an admitted,
unfrozen generation in an unretired session. A frozen generation still
accepts evidence, old votes, QCs, finality receipts, and L1 control updates.
A verified terminal receipt completes the session's application quota and
ends its regional production, even if successor selection is still pending.
Local terminal commitment alone does not retire a replica: it may still be
needed for others' recovery until a terminal receipt is available.
An effective forced closure retires the targeted session under
\cref{ass:cut}. Retired sessions retain their data and signing histories.

\begin{table}[htbp]\centering\small
\caption{Parent selection within an admitted generation.}
\label{tab:parent}
\begin{tabularx}{\linewidth}{@{}p{.32\linewidth}Y Y@{}}
\toprule
Position & Parent & Parent certificate\\
\midrule
First DATA, $t=j=0$ & $A_{r,g}$ & ANCHOR QC in $(r,g,0)$\\
First DATA, $t>0,j=0$ & Previous term's terminal DATA & Its proposal QC\\
Later DATA, $j>0$ & Position $j-1$ of this term & Its proposal QC\\
Only DRAIN, $g>0,t=j=0$ & Full-quota opening $A_{r,g}$ & ANCHOR QC in $(r,g,0)$\\
\bottomrule
\end{tabularx}
\end{table}

\section{Protocols}
\label{sec:protocols}
Protocols 1 and 2 patch v3.1's steady-state and recovery boxes. Red text
marks substantive changes from v3.1; the addition of the session index
$r$ to otherwise unchanged messages is implicit throughout. Protocol 3
is new and is shown in red. Explanations and service assumptions are
kept outside the execution steps.

\begin{protocol}{alg:steady}{Propose, vote, commit, and regular termination}
\begin{enumerate}[label=\textbf{\arabic*)},itemsep=5pt]
\item \textbf{Propose.} The authorized leader $\ell_{r,g,t}$ obtains the
parent and its QC from \cref{tab:parent}. \RR{If this is a DATA generation,
use the next position $j<m_{r,g,t}$. If it is a drain generation, use only
its one DRAIN position. A complete or retired session proposes nothing.}
Broadcast the complete proposal
\[
 P=\langle\tagg{PROPOSE},r,g,t,j,B,Q_{\rm parent}\rangle_{\ell_{r,g,t}}
\]
and process it locally under Vote. The parent need not be committed.

\item \textbf{Vote.} Process quit evidence first; return if it freezes
the current generation. Check completeness and validity. Request missing
data without voting; reject invalid proposals. If already voted for $P$,
stop. Otherwise forward $P$, broadcast
\RR{$\langle\tagg{VOTE},r,g,t,j,\Hash(B)\rangle_v$}, and start its
timer at the local acceptance time $\tau_v(P)$, expiring at
\RR{$\tau_v(P)+2\Delta_{i_r}$}. Process unseen earlier-term proposals
and attached history by the same rule. A duplicate never restarts a timer.

\item \textbf{Commit and export.} On timer expiry, commit $B$ and its
uncommitted ancestors. Normal term rotation preserves timers; Quit cancels
unfinished generation timers. A QC alone starts no timer.
\RR{Send FINAL endorsements for newly committed DATA blocks, including
the terminal DATA ancestor of a committed DRAIN. Relay assembled receipts
to peers and the cross-region infrastructure.}

\item \textbf{Terminal QC.} On a valid term-terminal proposal QC, process
its unseen history under Vote. \RR{If it certifies DATA strictly below
$H_r$}, and the generation remains unfrozen, send its complete bundle
once to $\ell_{r,g,t+1}$ and enter $(r,g,t+1)$ if not already there or
beyond it. Record first entry $e_v(r,g,t+1)$ and disable completed-term
checks. Process terminal-QC chains in order; duplicates never reset entry.
\RR{A QC for DATA at $H_r$, or for DRAIN, admits no next term or fresh
position. Continue pending timers and evidence processing.}

\item \RR{\textbf{Complete the region.} On a valid $\FC_r(T_r)$, adopt
its DATA prefix, relay the complete terminal bundle, and mark session $r$
complete. Stop its new proposals, votes, progress checks, unfinished
timers, and local recovery. Continue serving data and running Protocol 3.
An effective, authenticated forced closure instead retires $r$ at its
specified DATA anchor. Neither case withdraws a completed DATA commitment.}
\end{enumerate}
\end{protocol}

Acceptance occurs only after complete validation and the first vote.
The term-entry clock starts at admission, before any first proposal by
that term's leader. Reaching $H_r$ in a proposal or QC is not a cross-region
entry event. With no Quit, the terminal DATA's own timer is sufficient;
no subsequent block is required for that commitment.

\paragraph{Complete quit proofs.}
For an admitted $(r,g)$, a proof is either: (a) matching
$\langle\tagg{BLAME},r,g,t\rangle$ signatures of total weight
$>F_{i_r}$ for an admitted term; (b) two different admissible fresh
proposals signed by its leader at the same $(r,g,t,j)$; or (c) two
admissible NEWVIEWs signed by $\ell_{r,g,0}$ selecting different openings.
The proposals include their admission and quota witnesses. An excess
DATA block or unauthorized DRAIN is rejected, not treated as a second
admitted position. Different openings constitute equivocation even if
comparable; different certificates for the same opening do not.
Evidence from any admitted term can freeze its generation. Old-session
or old-generation evidence is archival and cannot freeze a current
successor. Future proofs wait for the corresponding authenticated entry.

\begin{protocol}{alg:recovery}{Blame, quit, and recovery with restricted drain}
\begin{enumerate}[label=\textbf{\arabic*)},itemsep=5pt]
\item \textbf{Blame.} At first term entry, schedule the first check at
\RR{$e_v(r,g,t)+6\Delta_{i_r}$}, even before NEWVIEW or a proposal arrives.
\RR{Once the opening is authenticated, use $q=m_{r,g,t}$ for a DATA term
and $q=1$ for a drain generation.} For $p=1,\ldots,q$, check at
\[
 e_v(r,g,t)+(2p+4)\Delta_{i_r}.
\]
If the term remains current and unfrozen and fewer than $p$ distinct
fresh positions have triggered local votes, broadcast
$\langle\tagg{BLAME},r,g,t\rangle_v$ once.
\RR{The one DRAIN vote counts; ANCHOR votes do not. Check an already-due
deadline when its opening becomes known. Do not require positions beyond $q$.}

\item \textbf{Quit.} On the first complete proof for current $(r,g)$,
freeze the generation: stop its progress checks, new proposals, votes,
and normal transitions; cancel all unfinished proposal timers.
Broadcast the proof and record quit time $a_v$ once. Retain commitments
and process old votes and QCs. Submit equivocation to L1 without awaiting
its finalization. \RR{Continue receiving finality receipts and force events.}

\item \textbf{Status and entry.} \RR{Unless the session has completed
or been retired}, at
\[
 e_v(r,g+1,0)=\max\{a_v+\Delta_{i_r},d(r,g+1)\},
\]
read $\tagg{AUTH}_{r,g+1}$. Fix the highest-ranked retained lock or valid
known QC from generations through $g$ as entry lock $Q_v$
\RR{(including the retained base justification $Q_r^0$)}.
Enter $(r,g+1,0)$, mark it unfrozen, start its progress clock, and send
$\langle\tagg{STATUS},r,g+1,0,Q_v,E\rangle_v$ with complete dependencies
to $\ell_{r,g+1,0}$. Process buffered proofs immediately. Later QCs update
retained evidence, not this fixed lock. \RR{Keep the original $A_r$,
$H_r$, and voting snapshot.}

\item \textbf{New-view.} The leader waits $2\Delta_{i_r}$ from its own
entry. If still current and unfrozen, choose the highest valid known QC
$Q^*$ from generations through $g$ \RR{or the base justification if no
ordinary QC is known}, and broadcast
\[
 \langle\tagg{NEWVIEW},r,g+1,0,A,Q^*,E\rangle_{\ell_{r,g+1,0}},
\]
where $Q^*$ certifies $A$, with complete dependencies. Choose one opening
for this generation. \RR{An opening below $H_r$ resumes DATA; an opening
at $H_r$, including a DRAIN descendant, admits exactly one fresh DRAIN.}

\item \textbf{First vote.} Process quit evidence, then validate NEWVIEW
and its dependencies. If current and unfrozen, no ANCHOR vote was cast,
and $\rank(Q^*)\ge\rank(Q_v)$, forward NEWVIEW and broadcast
$\langle\tagg{ANCHOR\mbox{-}VOTE},r,g+1,0,\Hash(A)\rangle_v$.
Start no timer and consume no fresh position. On an ANCHOR QC
\RR{of weight $>F_{i_r}$}, the leader resumes Protocol 1 at $j=0$.
\end{enumerate}
\end{protocol}

The entry-lock test governs ANCHOR voting only. A subsequent fresh
proposal backed by a valid ANCHOR QC is processed under Protocol 1 even
at a replica that declined the ANCHOR vote. A locally committed terminal
block does not justify stopping this recovery before a terminal receipt
or valid forced closure is available.

\subsection{Boundary witnesses and the bootstrap event}
\label{sec:boundary}
For an effective map $\chi$, let $\tagg{NextLive}(i,\chi)$ be the first
other live region after $i$ in $\mathcal L_R$, searching at most one cycle.
If none exists, its value is $\bot$. A boundary witness has the form
\[
 \mathcal W=(\mathit{mode},r,r+1,i',A,\FC(A),\lambda_s,\Xi),
 \qquad \mathit{mode}\in\{\tagg{REGULAR},\tagg{FORCE}\},
\]
where $\Xi$ authenticates the control order, effective map at selection
height $\lambda_s$, and any force entry and closure evidence. It also
binds the successor authority and configuration. The anchor's hash is
fixed by the block and finality receipt inside $\mathcal W$, not by the
multisignature height tag. The accepted witness seals that successor and
selection height. Certificate encodings do not change boundary identity.

For a regular boundary, $A=T_r$, its ancestry contains exactly $TB_r$
DATA blocks since $A_r$, and selection starts at the terminal block's
signed L1 height. For a forced boundary, $A$ is the finalized DATA block
at $h_*$ named by the activated entry, with
$h(A_r)\le h_*\le H_r$. If $h_*=h(A_r)$, reuse the incoming anchor and
its historical receipt. No new receipt from the failed region is required
for a prefix it inherited unchanged.
Forced selection uses $\lambda_s\ge\max\{a_j,\lambda(A)\}$ and replays
all due controls through $\lambda_s$ before choosing the successor.
For genesis, the initial L1 reference supplies $\lambda(A_0)$.

At an unsealed boundary, process due controls in their authenticated
order and apply an applicable force before regular quota completion at
the same control height. A sealed boundary is immutable. A force targeting
an already departed session cannot replace its successor; its status map
still applies in the ordered control history. A later live update does
not erase a latched force obligation or revive retired signing authority.
If $i'=\bot$, retain the anchor and advance the authenticated control
cursor; the first admissible later prefix with an eligible successor seals
the waiting boundary. A newer local L1 head cannot reselect it afterward.

\begin{definition}[Local cross-region bootstrap event]
\label{def:bootstrap}
For a replica $v\in V_{i'}$, $\tagg{Ready}_v(\mathcal W)$ occurs on its
first complete validation of an admissible sealed witness naming $i'$,
the anchor and missing application history, and the matching finalized
$\tagg{AUTH}_{r+1,0}$. REGULAR requires terminal finality and the regional
quota; FORCE additionally requires finalized L1 through its activation
height and the closure witness. The event is suppressed if this replica
already installed that successor or has authenticated a later session.
An incomplete boundary stays pending until its data and authority arrive.
\end{definition}
This is the precise event at which a node recognizes that its region may
start the next session. Merely receiving a leading block at height $H_r$,
its QC, an announced force entry, or a newer local status map is insufficient.

\paragraph{Terminal recovery example.}
Let $h(A_r)=100$, $TB_r=3$, and $K=2$. The first term proposes DATA
101--102 and its terminal QC admits the next term, which has only DATA
103. Its QC starts no next term. If its own timer survives, the receipt
for DATA 103 enables the regular handler. If Quit cancels the timers and
the next generation opens on certified DATA 103, its sole fresh proposal
is one DRAIN. That DRAIN's timer commits DATA 103 as an ancestor and
produces the same terminal DATA receipt. Another Quit may require a new
generation and one further DRAIN. In every case, the next region starts
with DATA 104, whose parent is DATA 103.

\paragraph{Forced bootstrap example.}
Suppose a session opened at DATA 100 with $TB_r=10$. A multisignature
entry names last available DATA height 106, marks its region off, and
explicitly schedules force at L1 height 1200. Prospective successors
pre-fetch DATA 106 and its receipt. Once L1 through 1200 is finalized and
the ordered closure witness is valid, the selected successor's Ready
event opens $(r+1,0,0)$ at DATA 106; its first fresh DATA is 107.
A mere announcement cannot start that session. If a valid final prefix
already includes DATA 107 under old authority, a cut at 106 does not meet
the stated service contract.

\begin{protocol}{alg:bootstrap}{Watch the stream and bootstrap the next region}
\color{rrRed}
\begin{enumerate}[label=\textbf{\arabic*)},itemsep=5pt]
\item \textbf{Watch and fetch.} Every prospective successor continuously
tracks the active session through the cross-region infrastructure.
Receive leading blocks and QCs, fetch missing ranges and finality receipts
from retained holders, and independently poll finalized L1 for authorities
and status entries. Validate and cache complete dependencies. Do not vote
or start regional progress clocks merely for speculative stream updates.

\item \textbf{Handle regular completion.} On a complete terminal DATA
bundle, verify its receipt and exactly $TB_r$ DATA positions. Starting
from its signed L1 reference, replay the authenticated boundary controls.
If no applicable force is pending, derive
$i'=\tagg{NextLive}(i_r,\chi_{\rm eff}(\lambda_s))$ and validate the
regular boundary witness. If none is eligible, keep this anchor and
resume selection on later ordered control updates.

\item \textbf{Handle activated force.} On finalized L1 through $a_j$,
process an explicit $\tagg{FORCE}(r,h_*)$ against its targeted active
session and latch it once. Verify the last-available-height and closure
witness, fetch $B_{h_*}$ with its historical finality receipt and ancestry,
and replay controls through $\lambda_s\ge\max\{a_j,\lambda(B_{h_*})\}$.
Derive the successor from that effective status map. Wait if evidence,
data, or an eligible region is missing. A stale target cannot reopen an
already sealed boundary. Neither handler waits for a fresh response from
the old region when the required evidence is retained elsewhere.

\item \textbf{Enter on Ready.} When $\tagg{Ready}_v(\mathcal W)$ fires
for this replica's region, relay the complete bundle within that region.
Install session $r+1$, its exact anchor $A_{r+1}=A$, configuration, quota,
and opening control height $\lambda_s$. Set its retained and entry lock
$Q_v$ to the fixed base justification $Q_{r+1}^0$. Enter $(r+1,0,0)$, record
$e_v(r+1,0,0)$ once, mark it unfrozen, and start the first progress check
before any leader message. Send STATUS containing the fixed base anchor
justification and $\mathcal W$ to $\ell_{r+1,0,0}$. Process buffered
current-generation evidence. A duplicate never resets entry, quota, or locks.

\item \textbf{Announce and certify the opening.} After $2\Delta_{i'}$
from its entry, the initial leader, if still current and unfrozen,
broadcasts NEWVIEW selecting exactly $A_{r+1}$, with
$Q^*=Q_{r+1}^0$ and $\mathcal W$.
Members validate the fixed anchor, authority, and entry proof, then forward
NEWVIEW and cast their one ANCHOR vote under Protocol 2's first-vote rule.
ANCHOR voting starts no timer. No highest-QC choice from the old region
may substitute another block for this opening.

\item \textbf{Start the region.} On the new region's ANCHOR QC of weight
$>F_{i'}$, its initial leader invokes Protocol 1 for the first DATA block
at global height $h(A_{r+1})+1$ and local physical position $1$.
Use the new session's fresh $TB_{r+1}$ quota. Subsequent local failures
invoke Protocol 2 within this session; they do not themselves increment $r$.
Continue watching L1 and streaming the new region's blocks and receipts.
\end{enumerate}
\end{protocol}

Both event handlers execute after every relevant stream or control update,
including while regional consensus is frozen. They are idempotent and use
the same ordered control state. Before installing or producing under a
successor, replay any further authenticated controls already due for that
successor. A force latches retirement even if it must wait for anchor data
or another live region; a later status recovery cannot cancel that latch.
Genesis initialization uses a synthetic initial boundary with the same
bootstrap exchange.

\paragraph{Entry clocks and incomplete bundles.}
The first local entry in Protocol 3 is the complete Ready event, not the
remote leader's proposal time, private QC assembly, or L1 installation
time. \Cref{ass:infra} supplies the destination entry skew needed for the
in-region progress checks. Cross-region transfer may take longer than
$\Delta_{i'}$ before this event. Once entry occurs, a silent initial leader
still faces the first check at $e_v(r+1,0,0)+6\Delta_{i'}$.

\begin{thebibliography}{9}
\bibitem{sync}
I. Abraham, D. Malkhi, K. Nayak, L. Ren, and M. Yin.
\emph{Sync HotStuff: Simple and Practical Synchronous State Machine Replication}.
IEEE Symposium on Security and Privacy, 2020, pp. 106--118.
\href{https://conferences.computer.org/sp/pdfs/sp/2020/349700a106.pdf}{Published paper}.
\end{thebibliography}
\end{document}
